\documentclass{article}
\usepackage{../assignment_preamble}

\title{Weekly Problem 5}
\author{Ravi Kini}
\date{May 4, 2026}

\begin{document}

\maketitle

\problem
\subproblem{(a)}
\begin{equation*}
\begin{nd}[rules=fde,justformat=just]
\hypo {1} {(p \lor q) \Rightarrow r}
\have {2} {\lnot (p \lor q) \lor r} \mce{1}
\open
\hypo {3} {\lnot (p \lor q)}
\have {4} {\lnot p} \nde{3}
\have {5} {\lnot\lnot r \lor \lnot p} \oi{4}
\close
\open
\hypo {6} {r}
\have {7} {\lnot\lnot r} \dn{6}
\have {8} {\lnot\lnot r \lor \lnot p} \oi{7}
\close
\have {9} {\lnot\lnot r \lor \lnot p} \oe{2,3-5,6-8}
\have {10} {\lnot r \Rightarrow \lnot p} \mci{9}
\end{nd}
\end{equation*}
We then see that $(p \lor q) \Rightarrow r \vdash \lnot r \Rightarrow \lnot p$ is a valid argument in the Logic of First-Degree Entailment.

The derivation only used rules that were available in both Strong Kleene logic and the Logic of Paradox, so the argument is valid in those logics as well.

\subproblem{(b)}
Let $p = b$ and $q = n$. We then see that $\lnot p = b$, so $\lnot p \Rightarrow q = 1$ and $\lnot p \Rightarrow q = 1$. However, $q \Rightarrow q = n$. Consequently, the argument $p, (p \land q) \Rightarrow r \vdash q \Rightarrow r$ is an invalid argument in the Logic of First-Degree Entailment.

However, this argument is valid in Strong Kleene logic and the Logic of Paradox, as seen by the following derivation.
\begin{equation*}
\begin{nd}[rules=strongkleene,justformat=just]
\hypo {1} {\lnot p \Rightarrow q}
\hypo {2} {p \Rightarrow q}
\have {3} {\lnot\lnot p \lor q} \mce{1}
\have {4} {\lnot p \lor q} \mce{2}
\open
\hypo {5} {\lnot p}
\have {6} {q} \ie{1,5}
\have {7} {\lnot q \lor q} \oi{6}
\close
\open
\hypo {8} {q}
\have {9} {\lnot q \lor q} \oi{8}
\close
\have {10} {\lnot q \lor q} \oe{3,5-7,8-9}
\have {11} {q \Rightarrow q} \mci{10}
\end{nd}
\end{equation*}
Th argument only fails in the Logic of First-Degree Entailment because of the truth values $b$ and $n$, which were not present in Strong Kleene logic or the Logic of Paradox.
\end{document}